% Minimal retained preamble for lane-p2 cards taken from
% sources/logic-fl02-conservation-of-ramsey-s-theorem-f/pinned-original-source.tex
% (Conservation of Ramsey's theorem for pairs, arXiv:2402.11616v3).
%
% Copied verbatim from the author's own file:
%   line 2      : \documentclass{article}
%   lines 16-25 : the author's theorem declarations
%   lines 46-88 : the author's macro definitions for the reverse-mathematics principles
%   line 91     : \card
%   lines 104-107: the natural sums and products, which rely on mathabx (source line 103)
%   line 116    : \qt
%
% Omitted author preamble lines and why (nothing was renamed or redefined):
%   line 3 (graphicx), 7 (url), 9 (cancel): no figure, url or cancel use in the retained spans;
%   line 6 (cleveref): the assembler adds its own cleveref machinery, a second load clashes;
%   line 8 (\usepackage[utf8x]{inputenc}): inputenc with the utf8x option fights the forced
%       xelatex engine;
%   lines 28-37 (the rep@theorem machinery and \newreptheorem{maintheorem}): the retained spans
%       use no rep-theorem environment;
%   lines 40-43 (xcolor and the three author comment macros): not used by the retained spans;
%   lines 93-99 (the LS2 font encoding block for \smashtimes based on stix): the host TeX Live has
%       no stix directory and the author himself never uses the result -- line 105 maps \ntimes to
%       \dottimes, and line 110 shows \smashtimes already commented out in the source.
% mathabx is retained because \dotplus, \dottimes and \dotdiv come from it.
\documentclass{article}
\usepackage{amsmath,amssymb,amsthm}
\usepackage{mathabx}
\usepackage{cancel}
\usepackage{cleveref}

\newtheorem{theorem}{Theorem}
\numberwithin{theorem}{section}
\newtheorem{maintheorem}[theorem]{Main Theorem}
\newtheorem{lemma}[theorem]{Lemma}
\newtheorem{question}[theorem]{Question}
\newtheorem{proposition}[theorem]{Proposition}
\newtheorem{statement}[theorem]{Statement}
\newtheorem{remark}[theorem]{Remark}
\newtheorem{definition}[theorem]{Definition}
\newtheorem{corollary}[theorem]{Corollary}

\newcommand{\RCA}{\mathsf{RCA}}
\newcommand{\WKL}{\mathsf{WKL}}
\newcommand{\SRT}{\mathsf{SRT}}
\newcommand{\EM}{\mathsf{EM}}
\newcommand{\SEM}{\mathsf{SEM}}
\newcommand{\ADS}{\mathsf{ADS}}
\newcommand{\SCAC}{\mathsf{SCAC}}
\newcommand{\SADS}{\mathsf{SADS}}
\newcommand{\COH}{\mathsf{COH}}
\newcommand{\ACA}{\mathsf{ACA}}
\newcommand{\qvdash}{\operatorname{{?}{\vdash}}}
\newcommand{\nqvdash}{\operatorname{{?}{\nvdash}}}
\newcommand{\BSig}{\mathsf{B}\Sigma^0}
\newcommand{\ISig}{\mathsf{I}\Sigma^0}
\newcommand{\LSig}{\mathsf{L}\Sigma^0}
\newcommand{\BPi}{\mathsf{B}\Pi^0}
\newcommand{\IPi}{\mathsf{I}\Pi^0}
\newcommand{\LPi}{\mathsf{L}\Pi^0}
\newcommand{\IDel}{\mathsf{I}\Delta^0}
\newcommand{\BCDe}{\mathsf{BC}\Delta^0}
\newcommand{\WF}{\mathsf{WF}}
\newcommand{\RT}{\mathsf{RT}}
\newcommand{\BME}{\mathsf{BME}}
\newcommand{\TT}{\mathsf{TT}}
\renewcommand{\O}{\mathcal{O}}
\renewcommand{\L}{\mathcal{L}}
\newcommand{\C}{\mathcal{C}}
\newcommand{\R}{\mathcal{R}}
\renewcommand{\S}{\mathcal{S}}
\newcommand{\T}{\mathcal{T}}
\renewcommand{\P}{\mathcal{P}}
\newcommand{\M}{\mathcal{M}}
\newcommand{\Nc}{\mathcal{N}}
\renewcommand{\P}{\mathcal{P}}
\newcommand{\Q}{\mathcal{Q}}
\newcommand{\NN}{\mathbb{N}}
\newcommand{\PP}{\mathbb{P}}
\newcommand{\uh}[0]{{\upharpoonright}}
\newcommand{\card}{\operatorname{card}}
\newcommand{\nplus}{\dotplus}
\newcommand{\ntimes}{\dottimes}
\newcommand{\nminus}{\dotdiv}
\newcommand{\nsum}{\dot{\sum}}
\def\qt#1{``#1''}
\title{An infinite stable tournament has a transitive subtournament preserving a prescribed ordinal well-foundedness}
\author{Quentin Le Houérou, Ludovic Levy Patey, Keita Yokoyama}
\date{arxiv:2402.11616v3; CC BY 4.0}
\begin{document}
\maketitle
\noindent\textit{Editorial note.} The counted unit is the article's own Corollary cor:sem-wf at source lines 1060-1066 together with the proof environment that begins on the next line (1067-1070); the pairing is the ordinary adjacent one.\par\medskip
\noindent\textit{Editorial note.} The retained context at source lines 457-459 and 760-767 is the statement of the two results that the counted proof invokes, namely the omega-extension theorem and the main proposition of the SEM section; their proofs are not retained and they are treated as prerequisites.\par\medskip
\noindent\textit{Editorial note.} The retained context at source lines 457-459 carries the author citation of Simpson and Smith inside the title of the theorem, and the author own bibliography file biblio.bib from the pinned e-print of arXiv:2402.11616v3 is attached to this task unedited and complete, with the bibtex backend, so that the citation resolves. The interval part of the pre-card gate cannot run for paper sources, so the public-interval comparison for this card was made directly against the public index: the only interval recorded there for this source is 522-524 and the retained intervals here are 457-459, 760-767, 1060-1066 and 1067-1070.\par\medskip
\noindent\textit{Editorial note.} The retained preamble reproduces verbatim the author's document class at source line 2, his theorem declarations at source lines 16-25, his macro definitions at source lines 46-88 and 91, and his natural-operation macros at source lines 104-107, and it retains mathabx and cleveref; the stix-dependent LS2 block and the inputenc load with the utf8x option are not retained.\par\medskip
\noindent\textit{Editorial note.} The article's conservation theorems and its earlier COH, WKL and SEM lemmas are not counted here and are not used as premises beyond the two retained statements.\par\medskip
\noindent\textit{Allowed standard prerequisites (editorial).}\begin{itemize}
\item The SEM forcing and its main proposition. For a countable model of RCA\_0 plus collection for Sigma\textasciicircum{}0\_2 plus well-foundedness of the supremum of the ordinals omega to the natural product with omega below alpha, topped by a set Y, and for every set P Turing above the jump of Y with the model of P satisfying RCA\_0\textasciicircum{}*, the forcing produces a set G which is an M-infinite transitive subtournament, computes the jump of G plus Y, and whose model satisfies RCA\_0 plus collection plus WF(alpha); the author's own statement is retained at source lines 760-767.. This is the statement the counted corollary specialises: the corollary drops the auxiliary set P by choosing it with the omega-extension theorem.
\item Omega-extension of a model of RCA\_0\textasciicircum{}* into a model of RCA\_0\textasciicircum{}* plus WKL. Every countable model of RCA\_0\textasciicircum{}* can be omega-extended into a model of RCA\_0\textasciicircum{}* plus weak Konig's lemma; retained with this card at source lines 457-459.. The proof uses it to manufacture the set P above the jump of Y that the main proposition requires.
\end{itemize}\medskip
\begin{theorem}[Simpson–Smith~\cite{simpson1986factorization}]\label[theorem]{thm:wkl-rca0s}
Every countable model of~$\RCA_0^*$ can be $\omega$-extended into a model of~$\RCA_0^* + \WKL$.
\end{theorem}

\begin{proposition}\label[proposition]{prop:sem-wf}
Consider a countable model $\M = (M,S) \models \RCA_0 + \BSig_2 + \WF(\sup_{\beta < \alpha} \omega^{\beta \times \omega})$ topped by a set $Y \in S$, where~$\alpha \in M$ is a non-zero ordinal $\leq \epsilon_0$. Then for every infinite stable tournament $R$ in $S$ and every set~$P \gg Y'$ such that $\M[P] \models \RCA_0^*$, there exists $G \subseteq M$ such that  
\begin{enumerate}
    \item $G$ is an $M$-infinite transitive subtournament ;
    \item $P \geq_T (G \oplus Y)'$ ;
    \item $\M[G] \models \RCA_0 + \BSig_2 + \WF(\alpha)$.
\end{enumerate}
\end{proposition}

\begin{corollary}\label[corollary]{cor:sem-wf}
Consider a countable model $\M = (M,S) \models \RCA_0 + \BSig_2 + \WF(\sup_{\beta < \alpha} \omega^{\beta \times \omega})$ topped by a set $Y \in S$, where~$\alpha \in M$ is an ordinal $\leq \epsilon_0$. Then for every infinite stable tournament $R$ in $S$, there exists $G \subseteq M$ such that  
\begin{enumerate}
    \item $G$ is an $M$-infinite transitive subtournament ;
    \item $\M[G] \models \RCA_0 + \BSig_2 + \WF(\alpha)$.
\end{enumerate}
\end{corollary}

\begin{proof}
By \Cref{thm:wkl-rca0s}, there exists a set~$P \subseteq M$ such that $P \gg Y'$ and $\M[P] \models \RCA_0^*$.
The conclusion follows from~\Cref{prop:sem-wf}.
\end{proof}

\bibliographystyle{plain}
\bibliography{retained-biblio}
\end{document}
